% 課題の提出要領（R＋Stanの回）。各章の\section{課題}の直後で\inputする
今回の課題は大問1つ，小問3つです。次の要領で提出してください。
\begin{description}
	\item[提出するもの] 分析を実行するRコード(\texttt{.R})と，モデルを記述したStanコード(\texttt{.stan})の2つです。
	\item[ファイル名] \texttt{chNN\_学籍番号\_氏名.R}，\texttt{chNN\_学籍番号\_氏名.stan}としてください。NNは回の番号です。
	\item[動くこと] 提出されたコードは，こちらの環境でそのまま実行して確認します。提出する前に，新しいセッションを開いて最初から最後まで通して実行し，エラーが出ないことを確かめてください。動かないものは未提出扱いになります。
	\item[答えの書き方] 小問で求めた数値は，指定された名前のオブジェクトに入れてください。小問1の答えは\verb|ans1|，小問2は\verb|ans2|，小問3は\verb|ans3|です。複数の数値を答えるときは，\verb|ans2 <- c(mu_eap = 9.7, mu_map = 9.9)|のように名前つきベクトルにまとめること。要素の名前はそれぞれの課題文に書いてあります。
	\item[数値は丸めない] オブジェクトに入れる数値は，\textbf{計算した値をそのまま}にしてください。\verb|round|や\verb|signif|で丸めると，正しく計算できていても不正解になることがあります。画面に表示するときに丸めるのは構いませんが，提出するオブジェクトの中身は丸めないこと。
\end{description}
わからないところがあれば，曜日別TAか小杉までメールで連絡し，指導を受けてください。
